\begin{python0}
from solutions import *; clear()
\end{python0}
\sectionthree{Exercise: C-string library}

Look at the chapter 18 Characters and C-strings. There are a bunch of
useful C-string functions in the C-string library. Implement your own
C-string library in files \texttt{mystring.h} and \texttt{mystring.cpp}.

\texttt{strlen(s)}

returns the string length of C-string s

\texttt{strcpy(s, t)}

copy string t to s

\texttt{strclear(s)}

after calling this function s becomes "". This is the same as strcpy(s,
"").

\texttt{strcat(s, t)}

concatenate t to s

\texttt{strcmp(s, t)}

returns 0 if C-strings s and t are the same strings

\texttt{strreplace(s, source, target)}

replace all occurrences of character source by character target. For
instance if s is "hello world", after calling strreplace(s,
'l', 'm'),
s becomes "hemmo wormd".

\texttt{strleftstrip(s)}

remove all whitespace characters (' ',
'\textbackslash n',
'\textbackslash t') on the left of s. For
instance if s is " abc ", after calling strleftstrip(s), s becomes "abc
".

\texttt{strrightstrip(s)}

remove all whitespace characters (' ',
'\textbackslash n',
'\textbackslash t') on the right of s. For
instance if s is " abc ", after calling strrightstrip(s), s becomes "
abc".

There are actually a lot more C-string functions that comes with your
C/C++ compiler.

